\documentclass[11pt,a4paper]{article}
\usepackage{fontspec,amsmath,amssymb,amsthm,microtype}
\setmainfont{lmroman10-regular.otf}[BoldFont=lmroman10-bold.otf,ItalicFont=lmroman10-italic.otf,BoldItalicFont=lmroman10-bolditalic.otf]
\usepackage{xcolor}
\usepackage[most]{tcolorbox}
\usepackage{mathrsfs}
\providecommand{\lt}{<}
\providecommand{\gt}{>}
\setmonofont{DejaVuSansMono.ttf}[BoldFont=DejaVuSansMono-Bold.ttf,ItalicFont=DejaVuSansMono-Oblique.ttf,BoldItalicFont=DejaVuSansMono-BoldOblique.ttf,Scale=0.86]
\newfontfamily\leanfallback{STIXTwoMath-Regular.otf}[Scale=0.86]
\definecolor{leanrule}{gray}{0.62}
\definecolor{leanback}{gray}{0.962}
\definecolor{leantitle}{gray}{0.22}
\definecolor{leanmuted}{gray}{0.38}
\newtcolorbox{leanbox}[1]{enhanced,breakable,colback=leanback,colframe=leanrule,
  boxrule=0.4pt,arc=0.6mm,left=2.4mm,right=2.4mm,top=1.8mm,bottom=1.8mm,
  before skip=0.8em,after skip=1em,fonttitle=\small\bfseries,coltitle=leantitle,
  colbacktitle=leanback,titlerule=0.3pt,title={#1}}
\newlength{\leancharwidth}
\newenvironment{leancode}{\par\smallskip\begingroup\ttfamily\footnotesize\raggedright
  \settowidth{\leancharwidth}{M}\setlength{\parskip}{0pt}\setlength{\parindent}{0pt}
  \frenchspacing}{\par\endgroup\smallskip}
\newcommand{\leanline}[2]{\hangindent=\dimexpr#1\leancharwidth+4\leancharwidth\relax
  \hangafter=1\noindent\hspace*{#1\leancharwidth}#2\par}
\newcommand{\leansource}[1]{\par{\footnotesize\color{leanmuted}#1}\par}
\newcommand{\leanlabel}[1]{\par\smallskip\noindent{\small\bfseries #1}\par}
\usepackage[margin=25mm]{geometry}
\usepackage[unicode,hidelinks]{hyperref}
\hypersetup{pdfauthor={LangevinLean Project}}
\setlength{\emergencystretch}{3em}
\setlength{\parskip}{0.4em}
\setlength{\parindent}{0pt}
\allowdisplaybreaks[2]
\raggedbottom
\title{\texorpdfstring{Actual shifted Langevin Euler convergence}{Actual shifted Langevin Euler convergence}}
\author{LangevinLean Project}
\date{}
\begin{document}\maketitle
Draft mathematical exposition; these notes may contain mistakes.


Subtracting the actual integral equation at two times yields the shifted equation started from the actual current state. Applying deterministic additive Euler convergence to this genuine \(L^2\) equation gives both fixed-terminal and moving-grid limits.


\section{\texorpdfstring{Restarting the actual integral equation}{Restarting the actual integral equation}}


\textbf{Claim.} For \(a,t\ge0\),

\[Y_{a+t}-Y_a=-\int_0^t\nabla V(Y_{a+s})\,ds+\sqrt2(B_{a+t}-B_a).\]

Euler recursion started from \(Y_a\) with these actual future increments converges in \(L^2\) to \(Y_{a+t}\), including along variable steps and counts whose elapsed time tends to \(t\).



\textbf{Proof.} Subtract the global integral equations at \(a+t\) and \(a\). Additivity and translation of the integrals give the shifted drift integral, while the exact Brownian increment theorem gives the shifted noise. The actual gradient map on \(L^2\) is Lipschitz and the global path is continuous, so the proved additive Euler convergence theorem applies.

\begin{leanbox}{Lean: \texttt{shifted\_\allowbreak{}state\_\allowbreak{}integral}, \texttt{shiftedEuler\_\allowbreak{}tendsto}, \texttt{shiftedEuler\_\allowbreak{}variable\_\allowbreak{}tendsto}}
\leansource{\texttt{shifted\_\allowbreak{}state\_\allowbreak{}integral} \quad WassersteinMixing/\allowbreak{}SemigroupEulerActual.lean, lines 20–33}
\begin{leancode}
\leanline{0}{theorem\ shifted\_state\_integral\ (P\ :\ Parameters)\ (V\ :\ E\ P.dimension\ →\ ℝ)}
\leanline{4}{(hV\ :\ PotentialAssumptions\ P\ V)\ (a\ T\ :\ ℝ)\ (ha\ :\ 0\ ≤\ a)\ (hT\ :\ 0\ ≤\ T)\ :}
\leanline{4}{letI\ :=\ actualProbability\_isProbability\ P\ V\ hV}
\leanline{4}{∀\ t\ ∈\ Icc\ 0\ T,\ state\ P\ V\ hV\ (a\ +\ t)\ -\ state\ P\ V\ hV\ a\ =}
\leanline{6}{-(∫\ s\ in\ 0..t,\ ULA.lpDrift\ (actualProbability\ P\ V)\ (gradient\ V)}
\leanline{8}{hV.gradientLipschitz\ (state\ P\ V\ hV\ (a\ +\ s)))\ +}
\leanline{6}{ULA.brownianLangevinNoise\ (actualProbability\ P\ V)\ (brownian\ P)}
\leanline{8}{(brownian\_preBrownian\ P\ V\ hV)\ a\ t}
\end{leancode}
\leansource{\texttt{shiftedEuler\_\allowbreak{}tendsto} \quad WassersteinMixing/\allowbreak{}SemigroupEulerActual.lean, lines 36–51}
\begin{leancode}
\leanline{0}{theorem\ shiftedEuler\_tendsto\ (P\ :\ Parameters)\ (V\ :\ E\ P.dimension\ →\ ℝ)}
\leanline{4}{(hV\ :\ PotentialAssumptions\ P\ V)\ (a\ t\ :\ ℝ)\ (ha\ :\ 0\ ≤\ a)\ (ht\ :\ 0\ ≤\ t)\ :}
\leanline{4}{letI\ :=\ actualProbability\_isProbability\ P\ V\ hV}
\leanline{4}{Tendsto\ (fun\ n\ :\ ℕ\ ↦\ ULA.additiveEulerChain}
\leanline{6}{(ULA.lpDrift\ (actualProbability\ P\ V)\ (gradient\ V)\ hV.gradientLipschitz)}
\leanline{6}{(t\ /\ ((n\ +\ 1\ :\ ℕ)\ :\ ℝ))\ (state\ P\ V\ hV\ a)}
\leanline{6}{(ULA.brownianLangevinNoise\ (actualProbability\ P\ V)\ (brownian\ P)}
\leanline{8}{(brownian\_preBrownian\ P\ V\ hV)\ a)\ (n\ +\ 1))\ atTop\ ({\leanfallback 𝓝}\ (state\ P\ V\ hV\ (a\ +\ t)))}
\end{leancode}
\leansource{\texttt{shiftedEuler\_\allowbreak{}variable\_\allowbreak{}tendsto} \quad WassersteinMixing/\allowbreak{}SemigroupEulerActual.lean, lines 55–76}
\begin{leancode}
\leanline{0}{theorem\ shiftedEuler\_variable\_tendsto\ \{I\ :\ Type*\}\ \{l\ :\ Filter\ I\}}
\leanline{4}{(P\ :\ Parameters)\ (V\ :\ E\ P.dimension\ →\ ℝ)\ (hV\ :\ PotentialAssumptions\ P\ V)}
\leanline{4}{(a\ T\ :\ ℝ)\ (ha\ :\ 0\ ≤\ a)\ (hT\ :\ 0\ ≤\ T)}
\leanline{4}{(step\ :\ I\ →\ ℝ)\ (count\ :\ I\ →\ ℕ)\ (hstep\ :\ ∀\ i,\ 0\ ≤\ step\ i)}
\leanline{4}{(hmesh\ :\ Tendsto\ step\ l\ ({\leanfallback 𝓝}\ 0))\ (hbound\ :\ ∀\ i,\ (count\ i\ :\ ℝ)\ *\ step\ i\ ≤\ T)}
\leanline{4}{(t\ :\ ℝ)\ (htime\ :\ Tendsto\ (fun\ i\ ↦\ (count\ i\ :\ ℝ)\ *\ step\ i)\ l\ ({\leanfallback 𝓝}\ t))\ :}
\leanline{4}{letI\ :=\ actualProbability\_isProbability\ P\ V\ hV}
\leanline{4}{Tendsto\ (fun\ i\ ↦\ ULA.additiveEulerChain}
\leanline{6}{(ULA.lpDrift\ (actualProbability\ P\ V)\ (gradient\ V)\ hV.gradientLipschitz)}
\leanline{6}{(step\ i)\ (state\ P\ V\ hV\ a)}
\leanline{6}{(ULA.brownianLangevinNoise\ (actualProbability\ P\ V)\ (brownian\ P)}
\leanline{8}{(brownian\_preBrownian\ P\ V\ hV)\ a)\ (count\ i))\ l\ ({\leanfallback 𝓝}\ (state\ P\ V\ hV\ (a\ +\ t)))}
\end{leancode}
\leanlabel{How the Lean proof follows the argument}
\texttt{shifted\_\allowbreak{}state\_\allowbreak{}integral} derives the actual restarted equation by interval-integral subtraction. \texttt{shifted\allowbreak{}Euler\_\allowbreak{}tendsto} uses equal elapsed subdivisions; \texttt{shifted\allowbreak{}Euler\_\allowbreak{}variable\_\allowbreak{}tendsto} applies the general moving-step theorem to the actual Lipschitz drift and shifted Brownian input.
\end{leanbox}


Earlier results used here: \cite[Section~3.13.2]{LangevinLean2026ULASC}; \cite[Section~3.13.5]{LangevinLean2026ULASC}; \cite[Section~3.16.1]{LangevinLean2026ULASC}; \cite[Section~3.16.9]{LangevinLean2026ULASC}; \cite[Section~3.25.1]{LangevinLean2026ULASC}; \cite[Section~3.25.3]{LangevinLean2026ULASC}; \cite[Section~3.44.2]{LangevinLean2026ULASC}.

\begin{thebibliography}{ULA-LSI}
\bibitem[ULA-SC]{LangevinLean2026ULASC}
LangevinLean Project. \emph{Unadjusted Langevin under Strong Convexity}. LL-01, 2026. Draft report, content revision d34ead06987c. \url{https://odakyildiz.github.io/LangevinLean/docs/pdf/UnadjustedLangevin/FullProof.pdf}.
\end{thebibliography}

\end{document}
